\clearpage
\section{The example system and its stationary states}\label{sec:stationary}

\noindent\motiv{\reve{The automaton of Example~\ref{ex:qca-rotation} rotates each oscillator of the
chain by a fixed angle and couples neighbouring oscillators through their momenta. It is the example
on which Kr\"uger demonstrates that a reversible Gaussian automaton can drive every finite region of
the chain towards equilibrium~\cite{krueger}, and this chapter prepares that result by describing the
automaton completely at each wave number and by determining its stationary covariances. At a single
wave number the automaton acts by a real $2\times2$ matrix of determinant one, and the analysis
rests on the elementary dichotomy for such matrices. When the trace lies strictly between $-2$ and
$2$, the matrix is conjugate to a rotation, its powers stay bounded, and it preserves exactly one
positive definite form up to a factor; when the absolute value of the trace exceeds two, it
squeezes, and the traces $\pm2$, which belong to $\pm1$ and to the parabolic matrices, separate the
two cases. Having determined the parameters for which every wave number is of the
first kind, the rotation angle as a function of the wave number, and the invariant form, we arrive
at a complete description of the stationary reflection-symmetric covariances: their symbols are the
invariant form multiplied, wave number by wave number, by a factor not less than one.}}

\medskip
\noindent\motiv{\reve{Throughout, $G=\Z$ and each site carries one mode, so that cells and symbols
are real or complex $2\times2$ matrices, and
$\Omega=\begin{psmallmatrix}0&1\\-1&0\end{psmallmatrix}$. Cell data, their covariance form and their
symbol are those of Definition~\ref{def:gs-cells}; for a rule $L$ and cell data $(\gamma_x)$ we
write $L_*\gamma$ for the cell data $(\gamma'_j)$ of Theorem~\ref{thm:gs-evolution}, whose covariance
form is $\gamma(L\cdot,L\cdot)$. Besides the identity $A^{\mathsf T}\Omega A=(\det A)\,\Omega$ of
Lemma~\ref{lem:qca-2x2}, one further identity for $2\times2$ matrices over a commutative ring is used
repeatedly,
\begin{equation}\label{eq:es-twobytwo}
\Omega A+A^{\mathsf T}\Omega=(\operatorname{tr}A)\,\Omega,
\end{equation}
which holds because, for $A=\begin{psmallmatrix}a&b\\c&d\end{psmallmatrix}$, the left side is
antisymmetric with upper right entry $a+d$; so is the Cayley--Hamilton relation
$A^2=(\operatorname{tr}A)A-(\det A)1$.}}

\subsection{Iterated evolution}

\noindent\motiv{\reve{The dynamics after $t$ steps is the automaton of the $t$-th power of the rule,
whose symbol is the $t$-th power of the symbol by Theorem~\ref{thm:qca-symbol}.}}

\begin{proposition}[Powers and iterates]\label{prop:es-iterate}
\reve{Let $L$, $L_1$, $L_2$ be rules and $(\gamma_x)$ cell data.
\begin{enumerate}[label=(\alph*)]
\item $\hat L(k+2\pi)=\hat L(k)$ and $\hat\gamma(k+2\pi)=\hat\gamma(k)$ for every $k\in\R$, and
$\widehat{L^t}(k)=\hat L(k)^t$ for every $t\in\mathbb N$.
\item $1_*\gamma=\gamma$ and $(L_1L_2)_*\gamma=L_{2*}(L_{1*}\gamma)$; hence $(L_*)^t\gamma=(L^t)_*\gamma$
for every $t\in\mathbb N$, and these are again cell data.
\item If $L$ is symplectic, so is $L^t$ for every $t$, and for a real matrix $S$ with
$S^{\mathsf T}\Omega S=\Omega$ the on-site rule $L_S$ is symplectic.
\end{enumerate}}
\end{proposition}

\begin{proof}\leavevmode
\reve{\begin{enumerate}[label=(\alph*)]
\item The evaluation points $e^{-\ii(k+2\pi)}$ and $e^{-\ii k}$ coincide, and the series for
$\hat\gamma(k+2\pi)$ and $\hat\gamma(k)$ agree term by term. Evaluation at $z=e^{-\ii k}$ is a ring
homomorphism (Lemma~\ref{lem:qca-circle}), and so is its entrywise application to matrices, which
therefore preserves powers.
\item Cell data are determined by their covariance form, because the form takes the value
$(\gamma_x)_{ij}$ on the site vectors $e_{x,i}$ and $e_{0,j}$. The rule $1$ is the identity, and
the form of $(L_1L_2)_*\gamma$ is $\gamma(L_1L_2\cdot,L_1L_2\cdot)$, which, since the rule of
$L_1L_2$ is $u\mapsto L_1(L_2u)$, is the form of $L_{1*}\gamma$ evaluated at $(L_2u,L_2w)$, that is,
the form of $L_{2*}(L_{1*}\gamma)$. The statement about iterates follows by induction on $t$ from
$(L_*)^{t+1}\gamma=L_*((L^t)_*\gamma)=(L^tL)_*\gamma$, and Theorem~\ref{thm:gs-evolution} shows at
each step that the result is cell data.
\item The symplectic symbols form a submonoid (Proposition~\ref{prop:gs-hom}), which contains the
powers of its elements. The involution fixes constant symbols and the constant embedding is a ring
homomorphism, so that the symplecticity condition for $L_S$ is the constant symbol of
$S^{\mathsf T}\Omega S=\Omega$.
\end{enumerate}}
\end{proof}

\noindent\motiv{\reve{Cell data can be recovered from their symbol by the inversion formula of
Fourier series, and two cell data with the same symbol therefore coincide; this is how
we will read off stationarity from the symbol.}}

\begin{proposition}[Fourier inversion]\label{prop:es-inversion}
\reve{Let $(\gamma_x)$ and $(\gamma'_x)$ be cell data.
\begin{enumerate}[label=(\alph*)]
\item For every $x\in\Z$, $\displaystyle\gamma_x=\frac1{2\pi}\int_{-\pi}^{\pi}e^{\ii kx}\,\hat\gamma(k)\,dk$,
entrywise.
\item If $\hat\gamma'(k)=\hat\gamma(k)$ for every $k$, then $\gamma'=\gamma$.
\end{enumerate}}
\end{proposition}

\begin{proof}\leavevmode
\reve{\begin{enumerate}[label=(\alph*)]
\item Each entry $\sum_g(\gamma_g)_{ab}e^{-\ii kg}$ of $\hat\gamma(k)$ is an absolutely convergent
Fourier series, which may be integrated term by term against $e^{\ii kx}$, and since
$\frac1{2\pi}\int_{-\pi}^\pi e^{\ii k(x-g)}\,dk=\delta_{g,x}$ only the term $g=x$ remains, as in
the proof of Lemma~\ref{lem:gs-parseval}.
\item By (a), both families are the same integrals of the same function.
\end{enumerate}}
\end{proof}

\begin{lemma}[Reflection-symmetric cells]\label{lem:es-refl}
\reve{Let $(\gamma_x)$ be cell data with $\gamma_{-x}=\gamma_x$ for every $x$. Then
$\hat\gamma(-k)=\hat\gamma(k)$, and $\hat\gamma(k)$ is a real symmetric matrix for every $k$.}
\end{lemma}

\begin{proof}
\reve{The substitution $x\mapsto-x$ gives
$\hat\gamma(-k)=\sum_x\gamma_{-x}e^{-\ii kx}=\hat\gamma(k)$, and by Lemma~\ref{lem:gs-cells}(b)
$\overline{\hat\gamma(k)}=\hat\gamma(-k)=\hat\gamma(k)$, so that every entry is real. A Hermitian
matrix with real entries is symmetric.}
\end{proof}

\subsection{The example system}

\noindent\motiv{\reve{We now fix the rule of Example~\ref{ex:qca-rotation}. Its coefficient at the
origin rotates each oscillator by an angle $\varphi$, and the coefficients at $\pm1$ couple each
site to its neighbours through the momentum row of that rotation, with strength $\lambda$.}}

\begin{definition}[The example system]\label{def:es-example}
\reve{For $\varphi,\lambda\in\R$, the \emph{example system} is the nearest-neighbour rule
$L=L_{\varphi,\lambda}$ of Example~\ref{ex:qca-rotation}, with coefficients
\[\ell_0=\begin{pmatrix}\cos\varphi&-\sin\varphi\\ \sin\varphi&\cos\varphi\end{pmatrix},\qquad
\ell_{1}=\ell_{-1}=\ell_\pm=\lambda\begin{pmatrix}\sin\varphi&\cos\varphi\\0&0\end{pmatrix},\]
and all other coefficients zero. We write
$\tau(k)=\cos\varphi+\lambda\cos k\sin\varphi$.}
\end{definition}

\begin{proposition}[The symbol]\label{prop:es-symbol}
\reve{For all $\varphi,\lambda,k\in\R$, the Fourier symbol of $L_{\varphi,\lambda}$ is the real
matrix
\[\hat L(k)=\ell_0+2\cos k\,\ell_\pm=
\begin{pmatrix}\cos\varphi+2\lambda\cos k\sin\varphi&-\sin\varphi+2\lambda\cos k\cos\varphi\\
\sin\varphi&\cos\varphi\end{pmatrix},\]
and
\begin{enumerate}[label=(\alph*)]
\item $\hat L(-k)=\hat L(k)$, and the rule is reflection symmetric, $\ell_{-g}=\ell_g$ for all $g$;
\item $\det\hat L(k)=1$, $\operatorname{tr}\hat L(k)=2\tau(k)$ and
$\hat L(k)^{\mathsf T}\Omega\hat L(k)=\Omega$;
\item the rule is symplectic;
\item for cell data $(\gamma_x)$, the evolved cell data have the symbol
$\widehat{L_*\gamma}(k)=\hat L(k)^{\mathsf T}\,\hat\gamma(k)\,\hat L(k)$.
\end{enumerate}}
\end{proposition}

\begin{proof}
\reve{The symbol is $\ell_\pm e^{\ii k}+\ell_0+\ell_\pm e^{-\ii k}=\ell_0+2\cos k\,\ell_\pm$, which is
real.
\begin{enumerate}[label=(\alph*)]
\item The cosine is even, and the only nonzero coefficients are $\ell_0$ and the equal coefficients
at $\pm1$.
\item The determinant is
$(\cos\varphi+2\lambda\cos k\sin\varphi)\cos\varphi-\sin\varphi(-\sin\varphi+2\lambda\cos k\cos\varphi)
=\cos^2\varphi+\sin^2\varphi$, we read off the trace, and the last identity is
Lemma~\ref{lem:qca-2x2} with $\det\hat L(k)=1$.
\item This is Example~\ref{ex:qca-rotation}.
\item This is Theorem~\ref{thm:gs-evolution}, whose reflection-symmetric form applies because
$L$ is reflection symmetric by (a); since $\hat L(k)$ is real, its transpose is its conjugate
transpose.
\end{enumerate}}
\end{proof}

\noindent\reve{In the regime the symbol is not symmetric, its off-diagonal entries being
$-\sin\varphi+2\lambda\cos k\cos\varphi$ and $\sin\varphi$, which would be equal only if
$\lambda\cos k\cos\varphi=\sin\varphi$, contrary to Lemma~\ref{lem:es-sign}; so the transpose in
(d) cannot be omitted.}

\begin{example}[The running example]\label{ex:es-running}
\reve{For $\varphi=\pi/2$ the on-site rotation is by a quarter turn, $\tau(k)=\lambda\cos k$, and
\[\hat L(k)=\begin{pmatrix}2\lambda\cos k&-1\\1&0\end{pmatrix}.\]
The rule maps the label $u(x)=(u_q,u_p)$ of a Weyl operator to
$\bigl(-u_p(x)+\lambda(u_q(x-1)+u_q(x+1)),\,u_q(x)\bigr)$; under the heuristic reading of
Proposition~\ref{prop:time-direction}, the field itself then moves by $L^{-1}$. For $\lambda=1$ the
rule is the inverse leapfrog of the massless lattice field, as the last section shows.}
\end{example}

\subsection{The regime of rotations}

\noindent\motiv{\reve{A real $2\times2$ matrix of determinant one with trace in $(-2,2)$ has the
eigenvalues $e^{\pm\ii\theta}$ with $\theta\in(0,\pi)$ and is conjugate to a rotation by $\theta$,
while a matrix of determinant one whose trace has absolute value greater than two has real
eigenvalues inverse to each other and stretches one direction while compressing another. Since
$\operatorname{tr}\hat L(k)=2\tau(k)$, every wave number is a rotation exactly when $|\tau|<1$
everywhere.}}

\begin{lemma}[Maximum of $|\tau|$]\label{lem:es-max}
\reve{For all $\varphi,\lambda,k$, $|\tau(k)|\le|\cos\varphi|+|\lambda\sin\varphi|$, with equality
for some $k$.}
\end{lemma}

\begin{proof}
\reve{The inequality is the triangle inequality with $|\cos k|\le1$. Equality holds at $k=0$ if
$\cos\varphi$ and $\lambda\sin\varphi$ are not of opposite signs, and at $k=\pi$ otherwise, since then
$\cos\varphi$ and $-\lambda\sin\varphi$ are not of opposite signs.}
\end{proof}

\noindent\motiv{\reve{Besides the rotation of every wave number, the convergence of Section~\ref{sec:convergence}
needs the chain to be coupled, since without coupling a squeezed initial state simply rotates for
ever at each site.}}

\begin{definition}[The regime]\label{def:es-regime}
\reve{The pair $(\varphi,\lambda)$ lies in the \emph{regime} if
\[|\cos\varphi|+|\lambda\sin\varphi|<1\qquad\text{and}\qquad\lambda\sin\varphi\neq0 .\]}
\end{definition}

\noindent\reve{The first condition already forces $\sin\varphi\neq0$ (Lemma~\ref{lem:es-sign} below),
so that the second only excludes the uncoupled chain $\lambda=0$.}

\begin{proposition}[Forms of the regime]\label{prop:es-regime}
\reve{Let $\varphi,\lambda\in\R$.
\begin{enumerate}[label=(\alph*)]
\item $|\cos\varphi|+|\lambda\sin\varphi|<1$ if and only if $|\operatorname{tr}\hat L(k)|<2$ for
every $k$.
\item If $\sin\varphi\neq0$, then $|\cos\varphi|+|\lambda\sin\varphi|<1$ if and only if
$|\lambda|<f_{\rm crit}$, where
\[f_{\rm crit}=\frac{1-|\cos\varphi|}{|\sin\varphi|}.\]
\item $f_{\rm crit}=\tan(\varphi/2)$ for $0<\varphi\le\pi/2$, and $f_{\rm crit}=\cot(\varphi/2)$ for
$\pi/2<\varphi<\pi$.
\end{enumerate}}
\end{proposition}

\begin{proof}\leavevmode
\reve{\begin{enumerate}[label=(\alph*)]
\item Since $|\operatorname{tr}\hat L(k)|=2|\tau(k)|$, Lemma~\ref{lem:es-max} shows that the maximum
over $k$ of $|\operatorname{tr}\hat L(k)|$ is $2(|\cos\varphi|+|\lambda\sin\varphi|)$ and is attained,
which is the claim.
\item Divide $|\lambda|\,|\sin\varphi|<1-|\cos\varphi|$ by $|\sin\varphi|>0$.
\item For $0<\varphi<\pi$ write $\varphi=2x$ with $\sin x,\cos x>0$, so that $\sin\varphi=2\sin
x\cos x>0$ while $1-\cos\varphi=2\sin^2x$ and $1+\cos\varphi=2\cos^2x$. If $\varphi\le\pi/2$, then
$\cos\varphi\ge0$ and $f_{\rm crit}=(1-\cos\varphi)/\sin\varphi=\sin x/\cos x$, and if
$\varphi>\pi/2$, then $\cos\varphi<0$ and $f_{\rm crit}=(1+\cos\varphi)/\sin\varphi=\cos x/\sin x$.
\end{enumerate}}
\end{proof}

\noindent\reve{For negative $\varphi$ the same computation, with $|\sin\varphi|=-\sin\varphi$,
gives $f_{\rm crit}=|\tan(\varphi/2)|$ for $-\pi/2\le\varphi<0$ and $|\cot(\varphi/2)|$ for
$-\pi<\varphi<-\pi/2$. In the running example $f_{\rm crit}=1$, and the regime is
$0<|\lambda|<1$.}

\noindent\motiv{\reve{Whether the invariant form of Section~\ref{sec:es-eps} is written with the
projector $P_k$ or with its conjugate depends on the sign of $\sin\varphi-\lambda\cos k\cos\varphi$
(Proposition~\ref{prop:es-eps}(d)), and the regime fixes this sign. The bound it needs goes beyond
the triangle inequality, because the regime bounds $|\lambda\sin\varphi|$ while the sign depends on
$|\lambda\cos\varphi|$.}}

\begin{lemma}[A sign]\label{lem:es-sign}
\reve{If $|\cos\varphi|+|\lambda\sin\varphi|<1$, then $\sin\varphi\neq0$,
$|\lambda\cos\varphi|<|\sin\varphi|$, and $\sin\varphi-\lambda\cos k\cos\varphi$ has the sign of
$\sin\varphi$ for every $k$.}
\end{lemma}

\begin{proof}
\reve{Write $a=|\cos\varphi|$ and $b=|\sin\varphi|$. Since $a<1$ and $a^2+b^2=1$, we have $b>0$.
Multiplying $|\lambda|b<1-a$ by $a\ge0$ gives $|\lambda|b\,a\le(1-a)a$, and since $1-a>0$ and
$a<1+a$, also $(1-a)a<(1-a)(1+a)=b^2$. Hence $|\lambda|a\,b<b^2$, and division by $b$ gives
$|\lambda\cos\varphi|<b$. Then $|\lambda\cos k\cos\varphi|\le|\lambda\cos\varphi|<|\sin\varphi|$, so
that $\sin\varphi-\lambda\cos k\cos\varphi$ differs from $\sin\varphi$ by less than
$|\sin\varphi|$ and therefore has the sign of $\sin\varphi$.}
\end{proof}

\subsection{The rotation angle}

\noindent\motiv{\reve{In the regime each $\hat L(k)$ is a rotation in suitable coordinates, and its
angle, which depends on $k$ unless $\lambda=0$, is determined by the trace.}}

\begin{definition}[Rotation angle]\label{def:es-alpha}
\reve{For $(\varphi,\lambda)$ in the regime,
$\theta(k)=\arccos\tau(k)=\arccos(\cos\varphi+\lambda\cos k\sin\varphi)$.}
\end{definition}

\begin{proposition}[The angle]\label{prop:es-alpha}
\reve{Let $(\varphi,\lambda)$ lie in the regime. Then $|\tau(k)|<1$, $\cos\theta(k)=\tau(k)$,
$\operatorname{tr}\hat L(k)=2\cos\theta(k)$, $\theta(k)\in(0,\pi)$ and $\sin\theta(k)>0$ for every
$k$, and
\begin{enumerate}[label=(\alph*)]
\item $\theta(-k)=\theta(k)$, and $\theta$ is continuous;
\item $\theta$ is differentiable, with $\theta'(k)=\lambda\sin\varphi\sin k/\sin\theta(k)$;
\item on $[0,\pi]$, $\theta$ is strictly increasing if $\lambda\sin\varphi>0$ and strictly decreasing
if $\lambda\sin\varphi<0$.
\end{enumerate}}
\end{proposition}

\begin{proof}
\reve{By Lemma~\ref{lem:es-max} and the regime, $|\tau(k)|<1$. On $[-1,1]$ the arccosine inverts the
cosine and maps $(-1,1)$ into $(0,\pi)$, where the sine is positive, and the trace is
Proposition~\ref{prop:es-symbol}(b).
\begin{enumerate}[label=(\alph*)]
\item $\tau$ is even and continuous, and the arccosine is continuous.
\item The arccosine has the derivative $-1/\sqrt{1-x^2}$ at $|x|<1$, and
$\sqrt{1-\tau(k)^2}=\sin\theta(k)$; with $\tau'(k)=-\lambda\sin\varphi\sin k$ the chain rule
gives the formula.
\item The cosine is strictly decreasing on $[0,\pi]$, so that $\tau$ is strictly decreasing there if
$\lambda\sin\varphi>0$ and strictly increasing if $\lambda\sin\varphi<0$, and its values lie in
$(-1,1)$, where the arccosine is strictly decreasing.
\end{enumerate}}
\end{proof}

\noindent\reve{Although (c) follows without derivatives, (b) adds, through $\lambda\sin\varphi\neq0$,
that $\theta'$ vanishes on $[0,\pi]$ only at the end points, so that $\theta$ is a change of
variables on each half of the circle, which is what the convergence argument will use. In the
running example, $\theta(k)=\arccos(\lambda\cos k)$ and $\sin\theta(k)=\sqrt{1-\lambda^2\cos^2k}$.}

\subsection{Powers of the symbol}

\noindent\motiv{\reve{After $t$ steps the symbol is $\hat L(k)^t$, and a matrix of determinant one
has powers of a simple closed form, since by the Cayley--Hamilton relation each power is a
combination of the matrix and the identity.}}

\begin{lemma}[Powers]\label{lem:es-chebyshev}
\reve{Let $A$ be a real $2\times2$ matrix with $\det A=1$ and $\operatorname{tr}A=2\cos a$, where
$\sin a\neq0$. Then for every $t\in\mathbb N$,
\[A^t=\frac{\sin(ta)}{\sin a}\,A-\frac{\sin((t-1)a)}{\sin a}\,1 .\]
In the regime this holds for $A=\hat L(k)$ and $a=\theta(k)$.}
\end{lemma}

\begin{proof}
\reve{Put $U_t=\sin(ta)/\sin a$, so that $U_0=0$ and $U_{-1}=-1$, and the formula
$A^t=U_tA-U_{t-1}1$ holds for $t=0$. If it holds for $t$, then by $A^2=2\cos a\,A-1$,
\[A^{t+1}=U_tA^2-U_{t-1}A=(2\cos a\,U_t-U_{t-1})A-U_t1,\]
and $2\cos a\,U_t-U_{t-1}=U_{t+1}$ is the identity $\sin((t+1)a)+\sin((t-1)a)=2\sin(ta)\cos a$. The
last statement follows from Propositions~\ref{prop:es-symbol}(b) and~\ref{prop:es-alpha}.}
\end{proof}

\noindent\motiv{\reve{The same powers can be written through the spectral projectors of $\hat L(k)$,
in terms of which Proposition~\ref{prop:es-eps}(e) expresses the invariant form. They are built from the real matrix
\[N(k)=\frac{\cos\theta(k)\,1-\hat L(k)}{\sin\theta(k)},\]
defined in the regime, which squares to $-1$ because, by the Cayley--Hamilton relation,
$(\cos\theta\,1-\hat L)^2=\cos^2\theta\,1-2\cos\theta\,\hat L+\hat L^2=-\sin^2\theta\,1$. The letter
$N$ also denotes neighbourhoods in these notes; with the argument $k$ it always means this matrix.}}

\begin{proposition}[Spectral projectors]\label{prop:es-projectors}
\reve{Let $(\varphi,\lambda)$ lie in the regime, let $P_k=\frac12\bigl(1+\ii N(k)\bigr)$, a complex
$2\times2$ matrix, and let $\bar P_k$ be its complex conjugate. Then $P_k+\bar P_k=1$, and
\begin{enumerate}[label=(\alph*)]
\item $P_k^2=P_k$ and $P_k\bar P_k=\bar P_kP_k=0$;
\item $\hat L(k)=e^{\ii\theta(k)}P_k+e^{-\ii\theta(k)}\bar P_k$, and $\hat L(k)P_k=e^{\ii\theta(k)}P_k$;
\item $\hat L(k)^t=e^{\ii t\theta(k)}P_k+e^{-\ii t\theta(k)}\bar P_k$ for every $t\in\mathbb N$.
\end{enumerate}}
\end{proposition}

\begin{proof}
\reve{Since $N=N(k)$ is real, $\bar P_k=\frac12(1-\ii N)$, and $P_k+\bar P_k=1$ uses nothing else.
\begin{enumerate}[label=(\alph*)]
\item With $N^2=-1$, $P_k^2=\frac14(1+2\ii N-N^2)=\frac12(1+\ii N)=P_k$, and
$P_k\bar P_k=\bar P_kP_k=\frac14(1+N^2)=0$, the two factors being polynomials in $N$.
\item $e^{\ii\theta}P_k+e^{-\ii\theta}\bar P_k=\cos\theta\,1-\sin\theta\,N=\hat L(k)$ by the
definition of $N$, and multiplying on the right by $P_k$ gives the second identity by (a).
\item By induction on $t$, the cross terms of the product with (b) vanishing by (a).
\end{enumerate}}
\end{proof}

\noindent\reve{Expanding (c) gives $\hat L(k)^t=\cos(t\theta)1-\sin(t\theta)N(k)$, which is
Lemma~\ref{lem:es-chebyshev} again, and it is this real form that the convergence argument will use.
The projector $P_k$ is Hermitian exactly when $N(k)$ is antisymmetric, which in the regime happens
only at $\cos k=0$, where the diagonal of $\tau(k)1-\hat L(k)$, namely
$(-\lambda\cos k\sin\varphi,\lambda\cos k\sin\varphi)$, vanishes; at every other wave number the two
projectors are not orthogonal.}

\subsection{The invariant form}\label{sec:es-eps}

\noindent\motiv{\reve{A rotation preserves the identity form, and a matrix conjugate to a rotation
preserves the transported form. For $\hat L(k)$ this form can be written down directly: since
$\hat L^{\mathsf T}\Omega\hat L=\Omega$, the matrix $\Omega(a1-\hat L)$ is preserved for every scalar
$a$, and by \eqref{eq:es-twobytwo} the choice $a=\cos\theta(k)=\frac12\operatorname{tr}\hat L(k)$ makes
it symmetric. Normalizing it to determinant one and orienting it to be positive definite gives the form of
the following definition.}}

\begin{definition}[Invariant form]\label{def:es-eps}
\reve{For $(\varphi,\lambda)$ in the regime,
\[\hat\varepsilon(k)=-\frac{\operatorname{sgn}(\sin\varphi)}{\sin\theta(k)}\,
\Omega\bigl(\cos\theta(k)\,1-\hat L(k)\bigr)=-\operatorname{sgn}(\sin\varphi)\,\Omega N(k).\]}
\end{definition}

\noindent\reve{Since $\cos\theta=\tau$, the entries of $\hat L(k)$ give
\begin{equation}\label{eq:es-eps}
\hat\varepsilon(k)=\frac{\operatorname{sgn}(\sin\varphi)}{\sin\theta(k)}
\begin{pmatrix}\sin\varphi&-\lambda\cos k\sin\varphi\\-\lambda\cos k\sin\varphi&\sin\varphi-2\lambda\cos
k\cos\varphi\end{pmatrix},
\end{equation}
and
\begin{equation}\label{eq:es-sin2}
\sin^2\theta(k)=1-\tau(k)^2=\sin^2\varphi-2\lambda\cos k\cos\varphi\sin\varphi-\lambda^2\cos^2k\sin^2\varphi
\end{equation}
is the determinant of the matrix in \eqref{eq:es-eps}. Without the factor
$\operatorname{sgn}(\sin\varphi)$ the matrix would be negative definite for $\sin\varphi<0$.}

\begin{proposition}[Properties of the invariant form]\label{prop:es-eps}
\reve{Let $(\varphi,\lambda)$ lie in the regime. Then $\hat\varepsilon(-k)=\hat\varepsilon(k)$ and
$\hat L(k)^{\mathsf T}\hat\varepsilon(k)\hat L(k)=\hat\varepsilon(k)$ for every $k$, and
\begin{enumerate}[label=(\alph*)]
\item $\hat\varepsilon(k)$ is real symmetric, and $\hat\varepsilon$ is continuous;
\item $\det\hat\varepsilon(k)=1$ and $(\Omega\hat\varepsilon(k))^2=-1$;
\item $\hat\varepsilon(k)$ is positive definite, with $\hat\varepsilon(k)_{11}=|\sin\varphi|/\sin\theta(k)$;
\item $\operatorname{tr}(\ii\Omega P_k)=(\sin\varphi-\lambda\cos k\cos\varphi)/\sin\theta(k)$, a real
number with the sign of $\sin\varphi$;
\item $\hat\varepsilon(k)=\operatorname{sgn}(\sin\varphi)\,\ii\Omega(P_k-\bar P_k)$, and for
$\sin\varphi>0$, $\hat\varepsilon(k)+\ii\Omega=2\ii\Omega P_k$.
\end{enumerate}}
\end{proposition}

\begin{proof}
\reve{Evenness follows from that of $\theta$ and $\hat L$. For the invariance, with $a=\cos\theta(k)$
and $\hat L=\hat L(k)$, the identity $\hat L^{\mathsf T}\Omega\hat L=\Omega$ gives
$\hat L^{\mathsf T}\Omega(a1-\hat L)\hat L=a\Omega-\Omega\hat L=\Omega(a1-\hat L)$, and
$\hat\varepsilon(k)$ is a scalar multiple of $\Omega(a1-\hat L)$.
\begin{enumerate}[label=(\alph*)]
\item Symmetry is visible in \eqref{eq:es-eps}, and continuity holds because $\theta$ and $\hat L$
are continuous and $\sin\theta$ does not vanish.
\item By \eqref{eq:es-eps} and \eqref{eq:es-sin2},
$\det\hat\varepsilon=\operatorname{sgn}(\sin\varphi)^2\sin^2\theta/\sin^2\theta=1$, since
$\sin\varphi\neq0$ by Lemma~\ref{lem:es-sign}. Further,
$\Omega\hat\varepsilon=-\operatorname{sgn}(\sin\varphi)\Omega^2N=\operatorname{sgn}(\sin\varphi)N$,
whose square is $N^2=-1$.
\item The $(1,1)$ entry is $\operatorname{sgn}(\sin\varphi)\sin\varphi/\sin\theta=|\sin\varphi|/\sin
\theta>0$, and a real symmetric $2\times2$ matrix $E$ with $E_{11}>0$ and $\det E>0$ is positive
definite, because $E_{11}\,v^{\mathsf T}Ev=(E_{11}v_1+E_{12}v_2)^2+(\det E)v_2^2$.
\item Since $\ii\Omega P_k=\frac\ii2\Omega-\frac12\Omega N$ and $\operatorname{tr}\Omega=0$, the
trace equals $-\frac12\operatorname{tr}(\Omega N)$, and the trace of $\Omega(\tau1-\hat L)$ is
$-2\sin\varphi+2\lambda\cos k\cos\varphi$; the sign is that of Lemma~\ref{lem:es-sign}.
\item Since $P_k-\bar P_k=\ii N$, we have $\ii\Omega(P_k-\bar P_k)=-\Omega N$, and for
$\sin\varphi>0$ the relation $\bar P_k=1-P_k$ gives
$\hat\varepsilon+\ii\Omega=\ii\Omega(P_k-\bar P_k+1)=2\ii\Omega P_k$.
\end{enumerate}}
\end{proof}

\noindent\reve{In the running example,
\[\hat\varepsilon(k)=\frac1{\sqrt{1-\lambda^2\cos^2k}}\begin{pmatrix}1&-\lambda\cos k\\-\lambda\cos
k&1\end{pmatrix}.\]
The entries of $\hat\varepsilon$ are real-analytic, $2\pi$-periodic
functions of $k$, since $1-\tau(k)^2$ is a trigonometric polynomial that does not vanish; their Fourier
coefficients therefore decay exponentially, and $\hat\varepsilon$ is the symbol of reflection-symmetric
cell data, as is $\eta\hat\varepsilon$ for every real-analytic, even, $2\pi$-periodic $\eta$. By
Corollary~\ref{cor:gs-purity}, (b) then says that these cell data saturate the state condition at
every wave number.}

\noindent\motiv{\reve{A factor multiplying $\hat\varepsilon(k)$ adds noise to the mode $k$, and the
state condition becomes a threshold on the factor.}}

\begin{proposition}[Threshold]\label{prop:es-threshold}
\reve{Let $E$ be a real symmetric positive definite $2\times2$ matrix with $\det E=1$, and $\eta\in\R$.
Then $\eta E+\ii\Omega\ge0$ if and only if $\eta\ge1$. In particular, in the regime,
$\eta\,\hat\varepsilon(k)+\ii\Omega\ge0$ if and only if $\eta\ge1$.}
\end{proposition}

\begin{proof}
\reve{Write $e_1=E_{11}>0$ and $e_2=E_{12}$. If $\eta E+\ii\Omega\ge0$, its $(1,1)$ entry $\eta e_1$
is non-negative, so that $\eta\ge0$, and its determinant, which by \eqref{eq:gs-detpm} is
$\eta^2\det E-1=\eta^2-1$, is non-negative; hence $\eta\ge1$. Conversely, for $w_\pm=\bigl(\sqrt{e_1},
(e_2\mp\ii)/\sqrt{e_1}\bigr)$ a direct computation, using $e_1E_{22}-e_2^2=1$, gives
$E\pm\ii\Omega=w_\pm w_\pm^\dagger$, and
\[\eta E+\ii\Omega=\tfrac{\eta+1}2(E+\ii\Omega)+\tfrac{\eta-1}2(E-\ii\Omega)\]
is a combination of positive semidefinite matrices with non-negative coefficients when $\eta\ge1$.
The last statement is the case $E=\hat\varepsilon(k)$, by Proposition~\ref{prop:es-eps}.}
\end{proof}

\subsection{Invariant forms of a rotation}

\noindent\motiv{\reve{To show that every reflection-symmetric stationary covariance is built from
$\hat\varepsilon$, we need that $\hat\varepsilon(k)$ spans the symmetric forms preserved by
$\hat L(k)$, and this follows from the description of the real $2\times2$ matrices commuting with a
given one.}}

\begin{lemma}[Commutant]\label{lem:es-commutant}
\reve{Let $A$ be a real $2\times2$ matrix that is not a multiple of the identity. Every real matrix
$X$ with $XA=AX$ is of the form $X=p1+qA$ with $p,q\in\R$.}
\end{lemma}

\begin{proof}
\reve{Subtracting multiples of the identity, we may take
$A_0=\begin{psmallmatrix}a&b\\c&-a\end{psmallmatrix}\neq0$ and
$X_0=\begin{psmallmatrix}x&y\\z&-x\end{psmallmatrix}$ with $X_0A_0=A_0X_0$. The commutator
$X_0A_0-A_0X_0$ has the entries $yc-bz$, $2(xb-ya)$, $2(za-xc)$ and $bz-yc$, so that
$(x,y,z)$ and $(a,b,c)$ have vanishing cross product, and $(x,y,z)=q(a,b,c)$ because
$(a,b,c)\neq0$.}
\end{proof}

\begin{proposition}[Invariant forms]\label{prop:es-invariant}
\reve{Let $A$ be a real $2\times2$ matrix with $\det A=1$ that is not a multiple of the identity.
\begin{enumerate}[label=(\alph*)]
\item A real symmetric matrix $S$ with $A^{\mathsf T}SA=S$ is of the form
$S=q\,\Omega\bigl(A-\tfrac12(\operatorname{tr}A)1\bigr)$ with $q\in\R$.
\item If $|\operatorname{tr}A|<2$ and $E\neq0$ is a real symmetric matrix with $A^{\mathsf T}EA=E$,
then every such $S$ is $S=\eta E$ for a unique $\eta\in\R$.
\end{enumerate}}
\end{proposition}

\begin{proof}\leavevmode
\reve{\begin{enumerate}[label=(\alph*)]
\item By Lemma~\ref{lem:qca-2x2}, $A^{\mathsf T}\Omega=\Omega A^{-1}$, so that $A^{\mathsf T}SA=S$
says that $X=\Omega^{-1}S$ commutes with $A$. By Lemma~\ref{lem:es-commutant}, $S=\Omega(p1+qA)$.
Symmetry of $S$, with $\Omega^{\mathsf T}=-\Omega$ and \eqref{eq:es-twobytwo}, becomes
$(2p+q\operatorname{tr}A)\Omega=0$, so that $p=-\frac12q\operatorname{tr}A$.
\item By (a), both $S$ and $E$ are multiples of $\Omega(A-\frac12\operatorname{tr}A\,1)$, the multiple
for $E$ being nonzero, and uniqueness holds because $E\neq0$.
\end{enumerate}}
\end{proof}

\noindent\reve{Conversely, $\Omega(A-\frac12\operatorname{tr}A\,1)$ is
symmetric and invariant by \eqref{eq:es-twobytwo}, and its determinant is
$1-\frac14(\operatorname{tr}A)^2$, positive exactly in the elliptic case. If the forms are not required
to be symmetric, the invariant ones span a plane, which also contains $\Omega$. This is why
reflection symmetry, which makes symbols real, enters the theorem below: the cell data with symbol
$2\hat\varepsilon(k)+\ii\sin k\,\Omega$ are stationary and satisfy the state condition, since
$2\hat\varepsilon(k)+\ii s\Omega=s\bigl((2/s)\hat\varepsilon(k)+\ii\Omega\bigr)\ge0$ for
$s=1+\sin k>0$ by Proposition~\ref{prop:es-threshold}, and trivially for $s=0$, but their symbol
is not of the form $\eta\hat\varepsilon$.}

\begin{corollary}[Coefficient]\label{cor:es-coeff}
\reve{In the regime, put $h_k(S)=-\frac12\operatorname{tr}\bigl(\Omega\hat\varepsilon(k)\Omega S\bigr)$
for a real matrix $S$. Then $h_k(\eta\,\hat\varepsilon(k))=\eta$, and every real symmetric $S$ with
$\hat L(k)^{\mathsf T}S\hat L(k)=S$ equals $h_k(S)\,\hat\varepsilon(k)$.}
\end{corollary}

\begin{proof}
\reve{By Proposition~\ref{prop:es-eps}(b),
$h_k(\eta\hat\varepsilon)=-\frac\eta2\operatorname{tr}\bigl((\Omega\hat\varepsilon)^2\bigr)=
-\frac\eta2\operatorname{tr}(-1)=\eta$. Since $|\operatorname{tr}\hat L(k)|<2$ and
$\hat\varepsilon(k)\neq0$ is symmetric and invariant, Proposition~\ref{prop:es-invariant}(b) gives
$S=\eta\hat\varepsilon(k)$ for some $\eta$, and then $\eta=h_k(S)$.}
\end{proof}

\subsection{The stationary states}

\noindent\motiv{\reve{A Gaussian state $(\omega,\gamma,d)$ is stationary under the automaton when
its covariance and its displacement are both preserved, $\gamma(L\cdot,L\cdot)=\gamma$ and $d\circ
L=d$. The theorem concerns the covariance; the displacement is settled in the remark after it. By
the preceding results, reflection-symmetric cell data are preserved exactly when their symbol is
preserved at every wave number, a preserved real symbol is a multiple of $\hat\varepsilon(k)$ at each
$k$, and the state condition then says that the multiple is at least one.}}

\begin{theorem}[Stationary covariances]\label{thm:es-stationary}
\reve{Let $(\varphi,\lambda)$ lie in the regime, and let $(\gamma_x)$ be cell data with
$\gamma_{-x}=\gamma_x$ for every $x$. The following are equivalent:
\begin{enumerate}[label=(\roman*)]
\item $L_*\gamma=\gamma$, and $\hat\gamma(k)+\ii\Omega\ge0$ for every $k$;
\item there is a continuous function $\eta\colon\R\to\R$ with $\eta(-k)=\eta(k)\ge1$ and
$\hat\gamma(k)=\eta(k)\,\hat\varepsilon(k)$ for every $k$.
\end{enumerate}}
\end{theorem}

\begin{proof}
\reve{(ii)$\Rightarrow$(i). By Proposition~\ref{prop:es-symbol}(d) and the invariance of
$\hat\varepsilon$, the evolved cell data have the symbol
$\hat L^{\mathsf T}\eta\hat\varepsilon\hat L=\eta\hat\varepsilon=\hat\gamma$, and by
Proposition~\ref{prop:es-inversion}(b) they equal $\gamma$. The state condition is
Proposition~\ref{prop:es-threshold}.}

\reve{(i)$\Rightarrow$(ii). By Lemma~\ref{lem:es-refl}, $S(k)=\hat\gamma(k)$ is real symmetric, even
in $k$, and continuous by Lemma~\ref{lem:gs-cells}. Taking symbols in $L_*\gamma=\gamma$ gives
$\hat L(k)^{\mathsf T}S(k)\hat L(k)=S(k)$, so that $S(k)=\eta(k)\hat\varepsilon(k)$ with
$\eta(k)=h_k(S(k))$ by Corollary~\ref{cor:es-coeff}. This $\eta$ is continuous and even, being a
trace of products of continuous even functions, and $\eta(k)\ge1$ by
Proposition~\ref{prop:es-threshold}.}
\end{proof}

\noindent\revs{The continuity and evenness of $\eta$ required in (ii), and the reflection symmetry
of the cells, are automatic:} \reve{if $\hat\gamma=\eta\hat\varepsilon$ for cell data
with $\eta\ge1$, then $\eta=\hat\gamma_{11}\sin\theta/|\sin\varphi|$ is continuous, the symbol is real,
and $\hat\gamma(-k)=\overline{\hat\gamma(k)}=\hat\gamma(k)$ makes $\eta$ even and, by
Proposition~\ref{prop:es-inversion}(b), the cells reflection symmetric. A stationary covariance of this
form has $\det\hat\gamma(k)=\eta(k)^2$. For $\eta\equiv1$ this is one at every wave number, the
saturated covariance, while $\eta>1$ adds noise to the corresponding modes; $\eta(k)$ measures the
excess noise of the mode $k$, as $\coth(\beta\omega/2)$ does for a thermal oscillator. For the
displacement, a translation-invariant $d(u)=\sum_x\delta^{\mathsf T}u(x)$ satisfies $d\circ L=d$ exactly
when $\hat L(0)^{\mathsf T}\delta=\delta$, and since $\hat L(0)$ is conjugate to a rotation by
$\theta(0)\in(0,\pi)$, it has no eigenvalue one, which forces $\delta=0$.}

\begin{example}[Stationary covariances of the running example]\label{ex:es-running-stat}
\reve{For $\varphi=\pi/2$ and $0<|\lambda|<1$, the stationary reflection-symmetric cell data satisfying
the Fourier criterion are exactly those with
\[\hat\gamma(k)=\frac{\eta(k)}{\sqrt{1-\lambda^2\cos^2k}}\begin{pmatrix}1&-\lambda\cos k\\-\lambda\cos
k&1\end{pmatrix},\qquad\eta(-k)=\eta(k)\ge1,\]
where $\eta$ is continuous and such that this is the symbol of cell data, for instance
real-analytic and $2\pi$-periodic.
The vacuum, $\hat\gamma\equiv1$, is reflection symmetric and satisfies the Fourier criterion, so that
by Theorem~\ref{thm:es-stationary} it would be stationary only if its symbol had this form; its
off-diagonal entry vanishes, so it is not stationary, and whether it approaches one of these states is
the question of convergence.}
\end{example}

\subsection{The massless field at the edge of the regime}

\noindent\motiv{\reve{The lattice Klein--Gordon field of Example~\ref{ex:gs-leapfrog} is a
nearest-neighbour rule, and one may ask whether every one of its wave numbers is a rotation. None of
its masses makes every wave number elliptic, and only the massless field is an example system, which
lies on the boundary of the regime.}}

\begin{proposition}[The leapfrog and the regime]\label{prop:es-finding2}
\reve{Let $c_k=2\cos k-m^2$ for a real $m^2$.
\begin{enumerate}[label=(\alph*)]
\item The symbols $\hat F$ and $\hat F^{-1}$ of the leapfrog and its inverse
(Example~\ref{ex:gs-leapfrog}) have $\operatorname{tr}\hat F(k)=\operatorname{tr}\hat F^{-1}(k)=c_k$,
and $|c_{k_0}|\ge2$ for $k_0=\pi$ if $m^2\ge0$ and $k_0=0$ if $m^2<0$.
\item The inverse leapfrog at $m^2=0$ is the example system with $\varphi=\pi/2$ and $\lambda=1$.
\item The pair $(\pi/2,1)$ is not in the regime: $|\cos\frac\pi2|+|\sin\frac\pi2|=1$, and
$\lambda=1=f_{\rm crit}(\pi/2)$.
\end{enumerate}}
\end{proposition}

\begin{proof}\leavevmode
\reve{\begin{enumerate}[label=(\alph*)]
\item Both traces are $c_k$ by Example~\ref{ex:gs-leapfrog}, and $|c_\pi|=2+m^2$ for $m^2\ge0$,
while $|c_0|=2-m^2>2$ for $m^2<0$.
\item At $m^2=0$ the inverse leapfrog has $\ell_0=\begin{psmallmatrix}0&-1\\1&0\end{psmallmatrix}$
and $\ell_{\pm1}=\begin{psmallmatrix}1&0\\0&0\end{psmallmatrix}$, which are the coefficients of
Definition~\ref{def:es-example} for $\varphi=\pi/2$ and $\lambda=1$.
\item $\cos\frac\pi2=0$ and $\sin\frac\pi2=1$.
\end{enumerate}}
\end{proof}

\noindent\reve{At $\lambda=1$ the running example has $\operatorname{tr}\hat L(0)=2$ and
$\operatorname{tr}\hat L(\pi)=-2$, and neither $\hat L(0)$ nor $\hat L(\pi)$ is $\pm1$, so that the
strict inequality of the regime fails exactly at the two parabolic modes $k=0$ and $k=\pi$. For
$m^2\neq0$, part (a) gives a mode with $|\operatorname{tr}|>2$, which squeezes.}

\subsection{Notes and references}\label{sec:es-notes}

\begin{enumerate}
\item \reve{\emph{The example system.} Definition~\ref{def:es-example} and
Proposition~\ref{prop:es-symbol} are (4.18) of~\cite{krueger}, where the coupling is written $f$
and the symbol $\hat\Gamma(k)=\Gamma_0+2f\cos k\,\Gamma_\pm$~\cite[p.~82]{krueger}. The evolved
symbol is (4.17b), printed with the conjugate transpose, and for this real symbol the conjugate
transpose and the transpose coincide.}
\item \reve{\emph{The regime.} The trace formula and the angle are (4.19), and $f_{\rm crit}$ with its
half-angle forms is (4.20)~\cite[p.~83]{krueger}, printed with absolute values, as here. The thesis
states the regime with $|f|\le f_{\rm crit}$ and restricts to the nondegenerate case, which we read as
the strict inequality of Definition~\ref{def:es-regime}~\cite[pp.~83--84]{krueger}. The thesis
excludes $\sin\varphi=0$ as degenerate~\cite[p.~88]{krueger} but not $\lambda=0$, and its
Proposition~4.7 does not hold as stated for $\lambda=0$, since an initial squeezed state then rotates
for ever. Lemma~\ref{lem:es-sign} makes precise the appeal to ``the condition on $f$ from (4.20)''
behind the sign statement of~\cite[p.~88]{krueger}. Figure~4.3 of the thesis~\cite[p.~84]{krueger}
($f=0.4$, $\varphi=\pi/4$) shows the angle increasing on $[0,\pi]$, in accordance with
Proposition~\ref{prop:es-alpha}(c).}
\item \reve{\emph{Projectors.} Proposition~\ref{prop:es-projectors} is Lemma~4.6 and (4.24)
of~\cite[pp.~84--85]{krueger}; the thesis proves idempotence through eigenvectors (its footnote~8),
here through $N(k)^2=-1$.}
\item \reve{\emph{Stationary states.} The thesis writes the invariant form as
$\hat\varepsilon(k)=\ii\sigma_s(P_k-\bar P_k)$ for $0<\varphi<\pi$ and
$\ii\sigma_s(\bar P_k-P_k)$ for $-\pi<\varphi<0$, with the real expression (4.28) for the first
case~\cite[pp.~87,~89]{krueger}; Definition~\ref{def:es-eps} is one formula for both, by
Proposition~\ref{prop:es-eps}(e). Theorem~\ref{thm:es-stationary} is Theorem~4.8 of the thesis,
with the factor written $g$ there and $\eta$ here, and the rotation angle written $\alpha(k)$ there
and $\theta(k)$ here, since $\alpha$ denotes automorphisms in these notes. The thesis proves that
every $g\hat\varepsilon$ with $g\ge1$ is stationary and that the limit of its dynamics has this form,
and it states that \emph{all} stationary reflection-symmetric states are of this form; the
completeness half, (i)$\Rightarrow$(ii), rests on Proposition~\ref{prop:es-invariant}, which is not in
the thesis. We take the state condition in its Fourier form; by Theorem~\ref{thm:gs-fourier} and
Remark~\ref{rem:gs-fourier-general} it is equivalent to the state condition of the covariance form.}
\item \reve{\emph{The massless field.} Proposition~\ref{prop:es-finding2} is not in the thesis.}
\end{enumerate}
